\section{Introduction and summary}
\label{sec:introduction}

Topology-changing transitions in Calabi--Yau compactifications admit a
particularly direct interpretation in M-theory. At a conifold degeneration,
M2-branes wrapped on exceptional curves become massless. Their charges are the
homology classes of those curves, and condensation along the Higgs branch
connects the resolved and deformed Calabi--Yau threefolds
\cite{Greene:1995hu,Katz:1996,Witten:1996qb,Mohaupt:2004de}. Contractions of
compact del Pezzo surfaces lead instead to interacting five-dimensional fixed
points \cite{Seiberg:1996,Morrison:1996pp,Douglas:1996,Intriligator:1997}. These two kinds of transition are usually analyzed
separately: the former in terms of charged hypermultiplets and the latter in
terms of the Coulomb and Higgs branches of a local superconformal field theory.

This paper studies compact transitions in which both types of sector occur at
once. We call a transition \emph{cooperative} when its global deformation
space contains a direction with nonzero components in several local sectors
but contains no nonzero direction supported on any one of those sectors
alone. The basic question is global. Suppose that a compact Calabi--Yau
threefold contains isolated singularities $p$ which, in their noncompact
geometric-engineering limits, define five-dimensional theories $\cT_p$. Each
local theory can possess one or more Higgs-branch directions. Which
combinations of these directions remain available when the singularities are
embedded in the same compact threefold and coupled to the vector fields of the
compactification?

There is no reason for the answer to factorize. A complex deformation of the
compact threefold restricts to a deformation of every singular germ, but the
image of the restriction map
\begin{equation}
  \TangentDef:=\operatorname{Ext}^1_X(\Omega_X^1,\cO_X)
  \longrightarrow \bigoplus_p T^1_{X,p},
  \qquad
  T^1_{X,p}:=\operatorname{Ext}^1_{\cO_{X,p}}
  (\Omega^1_{X,p},\cO_{X,p}),
  \label{eq:intro-restriction}
\end{equation}
can be a proper subspace. The same phenomenon has a five-dimensional
description. Reducing the M-theory three-form on harmonic two-forms gives
Abelian vector fields $A^I$. Wrapped M2 states are electrically charged under
these fields, while the compact embedding gauges Cartan subgroups of the
flavor symmetries of the interacting sectors. The corresponding moment-map
equations correlate local Higgs operators that would be independent in
isolated noncompact models.

We make this relation precise for compact Calabi--Yau threefolds whose isolated
singularities are ordinary double points and anticanonical cones over del
Pezzo surfaces of degrees seven and six. We write
\begin{equation}
 S_7=\operatorname{Bl}_2\mathbb P^2,
 \qquad
 S_6=\operatorname{Bl}_3\mathbb P^2 .
 \label{eq:intro-del-pezzo}
\end{equation}
The corresponding local sectors are a free hypermultiplet and the rank-one
$E_2$ and $E_3$ fixed points, respectively. This notation avoids the common
ambiguity between the degree of a del Pezzo surface and the number of blown-up
points.

The local input is finer than a dimension count. In one complex structure, a
nodal hypermultiplet with scalars $(h,\widetilde h)$ projects to the conifold
deformation coordinate $t=h\widetilde h$. For the $E_2$ theory, the invariant
ring of its Higgs chiral ring under the complexified $SU(2)$ Cartan is
\begin{equation}
 \CC[\cH_{E_2}]^{\CC^*}
 \cong \frac{\CC[J_0,U]}{(U^2,UJ_0)} .
 \label{eq:intro-e2-ring}
\end{equation}
This is the complete nonreduced miniversal base of the cone over $S_7$.
The reduced smoothing coordinate is the $SU(2)$ Cartan moment map $J_0$; the
additional nilpotent tangent direction is the $U(1)$ current-multiplet
operator $U$. For
$E_3$, taking invariants under the complexified maximal torus gives
\begin{equation}
 \CC[\cH_{E_3}]^{T_{\CC}}
 \cong
 \frac{\CC[H_1,H_2,L]}{(H_1L,H_2L)} .
 \label{eq:intro-e3-ring}
\end{equation}
The plane $(H_1,H_2)$ consists of the $SU(3)$ Cartan moment maps and the line
$L$ is the $SU(2)$ Cartan moment map. Equation
\eqref{eq:intro-e3-ring} is the reducible miniversal base of the cone over
$S_6$. Its plane is the maximal-torus invariant image
$\cB_{\mathfrak{sl}_3}\cong\CC^2$ of the minimal nilpotent orbit closure
$\cH_{\mathfrak{sl}_3}:=\overline{\mathcal O_{\min}(\mathfrak{sl}_3)}$;
its line is the corresponding image $\cB_{\mathfrak{sl}_2}\cong\CC$ of
$\cH_{\mathfrak{sl}_2}:=\overline{\mathcal O_{\min}(\mathfrak{sl}_2)}$.
This description keeps the two full $E_3$ Higgs components and their
invariant images separate rather than combining them into a single
tangent-space count. The local rings follow from
the known Higgs chiral rings
\cite{Cremonesi:2015lsa}; magnetic quivers provide an independent component
and Hilbert-series check \cite{Ferlito:2016,Cabrera:2018,Closset:2020afy}.

\subsection*{Charge constraints and the local-to-global deformation map}

Let $\wideX\to X$ be a crepant resolution and let $D_I$ be compact divisor
classes. At a node, denote the exceptional curve by $C_a$. At an $E_2$ or
$E_3$ point, let $R_{p\alpha}$ be a root in the summand of
$K_{S_p}^{\perp}$ selected by the chosen Higgs component. Here
$K_{S_p}^{\perp}$ is the orthogonal complement of the canonical class in
$H_2(S_p;\ZZ)$ for the intersection pairing. The compact matrix is
\begin{equation}
 c_{Ia}=D_I\cdot C_a,
 \qquad
 m_{Ip\alpha}=D_I\cdot\iota_{p*}R_{p\alpha}.
 \label{eq:intro-charges}
\end{equation}
The first set of entries are the electric charges of wrapped M2-branes. The
second set are the cocharacters specifying how $A^I$ gauges the local flavor
Cartans. Thus the complex moment maps take the form
\begin{equation}
 \mu_I^{\CC}
 =\sum_a c_{Ia}t_a
  +\sum_r m^{(7)}_{Ir}\beta_r
  +\sum_p\bigl(m_{Ip1}H_{p1}+m_{Ip2}H_{p2}+m_{Ip0}L_p\bigr),
 \label{eq:intro-moment-map}
\end{equation}
with $H_{p1}L_p=H_{p2}L_p=0$ at each $E_3$ point.

A \emph{branch profile} $\sigma$ chooses the full component
$\cH_{\mathfrak{sl}_2}$ or $\cH_{\mathfrak{sl}_3}$ at each $E_3$ point;
$\cB_{\sigma}$ is the product of their invariant images with the nodal and
reduced $E_2$ lines. Let $\cM_{H,\sigma}^{\mathrm{rigid}}$ denote the product
of the selected local Higgs cones, restricted to the triplet-moment-map zero locus
for the effective compact Abelian group and then divided by that group.
Restricting \eqref{eq:intro-moment-map} to the invariant coordinates gives a
linear complex map $\Qcompact$. For the geometric equality below, $X$ is a
$\Delta$-regular anticanonical hypersurface associated with an admissible
four-dimensional reflexive polytope as defined in
Section~\ref{subsec:admissible-class}, and it satisfies
\eqref{eq:global-hypotheses}. Our principal result is
\begin{equation}
 \boxed{
 \operatorname{Im}\!\left(
   \cM_{H,\sigma}^{\mathrm{rigid}}\longrightarrow\cB_{\sigma}
 \right)
 =\ker(\Qcompact)
 =\operatorname{Im}\!\left(
   \TangentDef\longrightarrow\bigoplus_pT^1_{X,p}
 \right)\cap\cB_{\sigma}.}
 \label{eq:intro-main}
\end{equation}
The first equality is a five-dimensional moment-map statement. The second is
the geometric local-to-global deformation theorem. Equation
\eqref{eq:intro-main} therefore identifies the geometric kernel with a physical
charge and gauging constraint, rather than attaching a field-theory
interpretation to a dimension computed after the fact.

Figure~\ref{fig:local-to-global} summarizes the three local inputs and the
common compact charge map.
\begin{figure}[t]
\centering
\begin{adjustbox}{max width=\textwidth,center}
\begin{tikzpicture}[
  node distance=7mm and 8mm,
  sector/.style={draw,rounded corners,fill=gray!7,align=center,
    minimum height=13mm,text width=0.25\textwidth,font=\small},
  map/.style={draw,very thick,rounded corners,align=center,
    minimum height=14mm,text width=0.22\textwidth,font=\small},
  result/.style={draw,double,rounded corners,align=center,
    minimum height=14mm,text width=0.22\textwidth,font=\small},
  arrow/.style={-{Latex[length=2.2mm]},thick}
]
\node[sector] (node) {nodes\\[-1mm]
  $t_a=h_a\widetilde h_a$\\
  wrapped-M2 hypers};
\node[sector,right=of node] (e2) {$E_2$ sectors\\[-1mm]
  $\beta_r=\mu^{\CC}_{SU(2)}$\\
  reduced Higgs line};
\node[sector,right=of e2] (e3) {$E_3$ sectors\\[-1mm]
  $(H_{p1},H_{p2})$ or $L_p$\\
  branch profile $\sigma$};

\node[map,below=12mm of e2] (charge) {compact charge map $Q_{X,\sigma}$\\[1mm]
  $c_{Ia}=D_I\!\cdot C_a$\,,\quad
  $m_{Ip\alpha}=D_I\!\cdot\iota_{p*}R_{p\alpha}$};
\node[result,right=15mm of charge] (kernel) {allowed holomorphic data\\[1mm]
  $\ker Q_{X,\sigma}$};

\draw[arrow] (node.south) -- ++(0,-5mm) -| ([xshift=-6mm]charge.north);
\draw[arrow] (e2.south) -- (charge.north);
\draw[arrow] (e3.south) -- ++(0,-5mm) -| ([xshift=6mm]charge.north);
\draw[arrow] (charge.east) -- node[above,font=\scriptsize]
  {$\mu_I^{\CC}=0$} (kernel.west);

\node[align=center,font=\scriptsize,below=5mm of charge,text width=0.66\textwidth]
  {The same divisor--curve pairing is the wrapped-M2 electric charge,
   the flavor-cocharacter gauging, and the dual of the geometric
   local-to-global map.};
\end{tikzpicture}
\end{adjustbox}
\caption{The compact Higgs-gluing mechanism for the sectors in the toric
examples, on a fixed $E_3$ branch profile.
The displayed kernel is the image in invariant complex moment-map
coordinates. Balanced representatives lift it to the zero level of the full
triplet moment maps. The exact geometric smoothing criterion requires this
kernel to meet the complement of every local discriminant.}
\label{fig:local-to-global}
\end{figure}


The word ``image'' in \eqref{eq:intro-main} is essential. A local deformation
base records invariant Higgs-branch coordinates, not the full
hyperk\"ahler cone or its metric. Likewise, the right-hand side is a complex
linear space on a fixed reduced profile. We do not identify it with the full
quaternionic-K\"ahler scalar manifold of the compact supergravity theory.

The real moment maps do not impose an additional inequality on this projected
image. Every value of $t_a$, $\beta_r$, $(H_{p1},H_{p2})$, or $L_p$ has a
balanced representative for which the corresponding local real Cartan moment
maps vanish. Consequently every point of the complex kernel lifts to the
zero-level triplet-moment-map locus. This also explains why signs in a homology
relation do not obstruct Higgsing: phases reside in the elementary complex
scalars, while their absolute values can be balanced.

The geometric comparison has an exact existence consequence. The mixed
smoothing theorem of \cite[Theorems~4.5 and~4.7]{Wuebben:2026smoothing},
combined with \eqref{eq:intro-main}, says that a convergent smoothing on a
chosen profile exists precisely when the charge kernel contains a vector
outside all local discriminants. At a node, an $E_2$ point or an $E_3$ line,
the local coefficient must be nonzero. On an $E_3$ plane, in the root
coordinates defined below, the condition is
\begin{equation}
 H_1H_2(H_1-H_2)\ne0.
 \label{eq:intro-e3-discriminant}
\end{equation}
A nonzero plane vector alone is insufficient. Necessity holds for arbitrary
analytic smoothing arcs, including those tangent to a discriminant to higher
order. If the relation space projects nontrivially to every selected $E_3$
summand, the selected global deformation base is smooth and every tangent
on it integrates. This hypothesis is automatic whenever the smoothing
criterion holds; nodes and $E_2$ points require no additional projection
hypothesis for smoothness of the selected base. These are geometric
integrability statements, not a construction of the compact scalar metric.

\subsection*{Three compact phenomena}

The theorem produces three qualitatively different effects.

First, a local Higgs direction can disappear completely after compact
embedding. We give a compact model $X_7$ with one $E_2$ cone point. The
reduced local theory has its $SU(2)$ Higgs component, but the compact kernel
is zero. Hence no global first-order deformation activates it and no analytic
smoothing exists.

Second, local sectors can Higgs only cooperatively. Our main explicit model $X_9$
contains two nodes and one $E_2$ point. Neither the two nodal sectors alone nor
the $E_2$ sector alone has a nonzero global direction. In a primitive compact
vector basis the effective charge matrix is
\begin{equation}
 Q_{X_9}^{\mathrm{eff}}=
 \begin{pmatrix}1&0&1\\0&1&1\end{pmatrix},
 \qquad
 \ker Q_{X_9}^{\mathrm{eff}}=\CC(-1,-1,1).
 \label{eq:intro-x9}
\end{equation}
The unique projected Higgs direction moves all three sectors together. An
explicit family integrates it. The Hodge numbers change as
\begin{equation}
 (h^{1,1},h^{2,1})=(6,122)\longrightarrow(3,123).
 \label{eq:intro-x9-hodge}
\end{equation}
Two vectors are Higgsed by the rank-two gauging, while the third disappearing
vector is the neutral rank-one $E_2$ Coulomb direction represented by the
exceptional divisor. The surviving deformation supplies one neutral
hypermultiplet, in agreement with \eqref{eq:intro-x9-hodge}.

Third, local smoothability need not imply simultaneous global Higgsability. A
compact model $X^{\circ}$ contains $26$ nodes, two $E_2$ points and two $E_3$
points. Every germ is locally smoothable. For all four ordered choices of
$\mathfrak{sl}_2$ and $\mathfrak{sl}_3$ components at the two $E_3$ points,
however, explicit primitive compact vectors force two nodal and
both $E_2$ deformation coordinates to vanish. The projected kernels have
dimensions $9,10,10,12$, so many partial Higgs directions survive, but no
profile contains a trajectory that activates all thirty sectors. The
vanishing-period theorem promotes this conclusion beyond tangent space on
every profile, because any smoothing must move the two forced nodes and both
$E_2$ points whatever components it selects at the $E_3$ points.
Here the superscript circle is part of the model name and does not denote the
polar operation introduced in \eqref{eq:polar-convention}.

\subsection*{Relation to previous work and scope}

Compact conifold charge matrices and their five-dimensional
gauged-supergravity actions were developed in \cite{Mohaupt:2004de}.
Malyshev's compact quintic supplied the closest geometric precursor: a del
Pezzo root curve is homologous to four conifold curves, reducing the number of
del Pezzo deformations \cite{Malyshev:2007zz}. Local $E_n$ Higgs branches and
magnetic quivers are treated in
\cite{Cremonesi:2015lsa,Ferlito:2016,Cabrera:2018,Bourget:2020bxh,Closset:2020afy}.
De Marco, Del Zotto, Grimminger and Sangiovanni develop the local Higgs-branch
stratification of five-dimensional fixed points from canonical singularities
\cite{DeMarco:2026higgs}. The present problem is the subsequent compact one:
which invariant directions survive when several such sectors share the
dynamical vectors of a single compact threefold. Compact gluing and multistage
decoupling are analyzed in
\cite{Apruzzi:2019vpe,Cvetic:2023plv}.

Our result extends the mixed geometry in
\cite{Malyshev:2007zz}. It gives a uniform theorem for free and interacting
localized sectors and derives its matrix as the electric-charge and
flavor-cocharacter gauging map. It also retains the two components of the
$E_3$ Higgs cone, identifies the full scheme-level $E_2/E_3$ projections, and
identifies the compact gauging with the relation map in the exact smoothing
criterion. The geometric necessity, sufficiency and selected-base
smoothness are results of \cite{Wuebben:2026smoothing}; the contribution here
is their combination with the local operator dictionaries and dynamical
compact gauging.
The cooperative and obstructed compact examples test each part of this
structure. None of these field-theoretic or branchwise statements is contained
in the one-example deformation count of \cite{Malyshev:2007zz}.

All compact harmonic modes used in the construction have finite norm at
finite volume, so the full matrix in \eqref{eq:intro-moment-map} is dynamical
in the compact theory. In a limit that decouples gravity, only vector
combinations with finite nonzero kinetic norm remain dynamical. The matrix is
then projected to $Q_{\mathrm{dyn}}=P_{\mathrm{dyn}}Q_X$. Gravity may therefore
decouple before all local sectors factorize; conversely, a fully factorized
local limit removes the shared rows. We do not claim that the compact
constraint survives every decoupling limit, and no example-specific K\"ahler
scaling ray is needed for the finite-volume theorem.

The paper is organized as follows. Section~\ref{sec:local-sectors} establishes
the local operator dictionaries and magnetic-quiver checks.
Section~\ref{sec:compact-gauging} derives the compact charge and gauging map.
Section~\ref{sec:global-theorems} states the first-order comparison, exact
smoothing criterion and selected-base integrability theorem.
Section~\ref{sec:examples} presents the absent, cooperative and obstructed
Higgsing examples, including an obstruction that forces an $E_3$ deformation
onto its discriminant. Section~\ref{sec:decoupling}
separates the compact result from its possible rigid limits, and
Section~\ref{sec:discussion} discusses extensions and open problems.
